\subsection*{3.4 Geometric Interpretation}

\begin{conceptbox}
    The directional derivative represents the slope of the surface in a given direction.
\end{conceptbox}

\begin{examplebox}
    For $f=x^2+y^2$, moving outward from origin increases value fastest.

    Directional derivative is highest along radial direction.
\end{examplebox}
Directional derivatives appear in:
\begin{itemize}
    \item heat transfer (temperature gradients)
    \item fluid flow (velocity change along direction)
    \item signal processing (gradient-based edge detection)
\end{itemize}

\begin{examplebox}
    Temperature field:
    \begin{align*}
        T(x,y)=x^2+y^2
    \end{align*}

    Directional derivative gives how fast temperature changes in a chosen direction.
\end{examplebox}
